% ========================================================================================
% CAPITOLO 2 - RETI SEQUENZIALI ELEMENTARI
% ========================================================================================
\chapter{Reti sequenziali elementari}

% ----------------------------------------------------------------------------------------
\section{Esercizio 3: Riconoscitore di sequenze}
% ----------------------------------------------------------------------------------------

\subsection{Esercizio 3.1}

\begin{traccia}
	Progettare, implementare in VHDL e testare mediante simulazione una macchina in grado di
	riconoscere la sequenza 101. La macchina prende in ingresso un segnale binario i che
	rappresenta il dato, un segnale A di tempificazione e un segnale M di modo, che ne
	disciplina il funzionamento, e fornisce un'uscita Y alta quando la sequenza viene
	riconosciuta. In particolare,
	\begin{itemize}[nosep]
		\item se M=0, la macchina valuta i bit seriali in ingresso a gruppi di 3 (sequenze non sovrapposte),
		\item se M=1, la macchina valuta i bit seriali in ingresso uno alla volta, tornando allo
		      stato iniziale ogni volta che la sequenza viene correttamente riconosciuta
		      (sequenze parzialmente sovrapposte).
	\end{itemize}
\end{traccia}

\progetto

La Tabella~\ref{tab:seq-esempio} chiarisce il comportamento atteso nei due modi su una
sequenza di prova.

\begin{table}[H]
	\centering\small
	\caption{Uscita \sig{Y} nei due modi per l'ingresso \texttt{11010101}.}
	\label{tab:seq-esempio}
	\begin{tabular}{@{}l*{8}{c}@{}}
		\toprule
		Posizione           & 1 & 2 & 3 & 4 & 5 & 6 & 7 & 8 \\
		\sig{i}             & 1 & 1 & 0 & 1 & 0 & 1 & 0 & 1 \\ \midrule
		$Y$ con \sig{M}=0   &   &   & 0 &   &   & \textbf{1} &   &   \\
		$Y$ con \sig{M}=1   & 0 & 0 & 0 & \textbf{1} & 0 & 0 & 0 & \textbf{1} \\
		\bottomrule
	\end{tabular}
\end{table}

\begin{scelta}[Modello di Moore]
	La macchina è realizzata secondo il modello di Moore: l'uscita dipende solo dallo stato
	e resta stabile per un intero periodo del segnale di tempificazione, così è ben visibile
	su un LED. \todo{Motivare, confrontare con Mealy.}
\end{scelta}

Nel modo \sig{M}=1 la macchina ha quattro stati, che rappresentano quanta parte della
sequenza è stata riconosciuta (Figura~\ref{fig:fsm-101}).

\begin{figure}[H]
	\centering
	\begin{tikzpicture}[node distance=2.6cm, auto, >={Stealth[length=2mm]}, every edge/.append style={thick}]
		\node[fsmstate, initial, initial text=] (S0) {$S_0$\\[-2pt]{\scriptsize $Y{=}0$}};
		\node[fsmstate, right=of S0] (S1) {$S_1$\\[-2pt]{\scriptsize $Y{=}0$}};
		\node[fsmstate, right=of S1] (S2) {$S_2$\\[-2pt]{\scriptsize $Y{=}0$}};
		\node[fsmstate, accepting, right=of S2, fill=osservcol!12] (S3) {$S_3$\\[-2pt]{\scriptsize $Y{=}1$}};
		\path[->]
			(S0) edge[loop above]              node {0} ()
			(S0) edge                          node {1} (S1)
			(S1) edge[loop above]              node {1} ()
			(S1) edge                          node {0} (S2)
			(S2) edge                          node {1} (S3)
			(S2) edge[bend left=40]            node {0} (S0)
			(S3) edge[bend left=50]            node {0} (S0)
			(S3) edge[bend right=40, swap]     node {1} (S1);
	\end{tikzpicture}
	\caption{Automa di Moore del riconoscitore nel modo \sig{M}=1. Da $S_3$ la macchina riparte
	come da $S_0$: i bit della sequenza riconosciuta non vengono riusati.}
	\label{fig:fsm-101}
\end{figure}

\todo{Automa per il modo \sig{M}=0 (conteggio della posizione nel gruppo di 3) e gestione del cambio di modo.}

\implementazione
\todo{Codice della FSM (tipo enumerato per gli stati, process di stato e process combinatorio).}

\simulazione
\todo{Testbench con la sequenza della Tabella~\ref{tab:seq-esempio} e casi limite.}

\subsection{Esercizio 3.2}

\begin{traccia}
	Sintetizzare e implementare su board la rete sviluppata al punto precedente, utilizzando
	uno switch S1 per codificare l'input i e uno switch S2 per codificare il modo M, in
	combinazione con due bottoni B1 e B2 utilizzati rispettivamente per acquisire l'input da
	S1 e S2 in sincronismo con il segnale di tempificazione A, che deve essere ottenuto a
	partire dal clock della board. Infine, l'uscita Y può essere codificata utilizzando un led.
\end{traccia}

\sintesiboard
\todo{Generazione di A dal clock (tick), acquisizione di i e M con B1/B2, vincoli, prova su scheda.}

% ----------------------------------------------------------------------------------------
\section{Esercizio 4: Shift register}
% ----------------------------------------------------------------------------------------

\subsection{Esercizio 4.1}

\begin{traccia}
	Progettare, implementare in VHDL e testare mediante simulazione un registro a scorrimento
	di N bit in grado di shiftare a destra o a sinistra di un numero Y variabile di posizioni
	a seconda di una opportuna selezione. In particolare, i valori possibili di Y sono 1 e 2.
	L'utente tramite selezione deve scegliere di quante posizioni shiftare. Il componente deve
	essere realizzato utilizzando sia un a) approccio comportamentale sia un b) approccio strutturale.

	Nota: il numero di bit del registro deve essere implementato come un generic, e
	dall'esterno deve poter essere scelta la modalità di funzionamento mediante opportuni
	segnali di selezione.
\end{traccia}

\progetto
\todo{Interfaccia (tabella \texttt{porte}), scelte sul riempimento dei bordi e sul caricamento parallelo.}

\implementazione
\todo{Architettura comportamentale e architettura strutturale (flip-flop + mux per ogni bit, \texttt{for generate}).}

\simulazione
\todo{Testbench che confronta le due architetture sugli stessi stimoli.}

% ----------------------------------------------------------------------------------------
\section{Esercizio 5: Cronometro}
% ----------------------------------------------------------------------------------------

\subsection{Esercizio 5.1}

\begin{traccia}
	Progettare, implementare in VHDL e testare mediante simulazione un cronometro, in grado di
	scandire secondi, minuti e ore a partire da una base dei tempi prefissata (es. si consideri
	il clock a disposizione sulla board). Il progetto deve prevedere la possibilità di
	inizializzare il cronometro con un valore iniziale, sempre espresso in termini di ore,
	minuti e secondi, mediante un opportuno ingresso di set, e deve prevedere un ingresso di
	reset per azzerare il tempo.

	Il componente deve essere realizzato utilizzando un approccio strutturale, collegando
	opportunamente dei contatori secondo uno schema a scelta.
\end{traccia}

\progetto

\begin{figure}[H]
	\centering
	\begin{tikzpicture}[node distance=1.6cm]
		\node[blocco] (tick) {base dei tempi\\\SI{1}{\hertz}};
		\node[blocco, right=of tick] (sec) {secondi\\mod 60};
		\node[blocco, right=of sec] (min) {minuti\\mod 60};
		\node[blocco, right=of min] (ore) {ore\\mod 24};
		\draw[seg] (tick) -- node[above, font=\scriptsize]{en} (sec);
		\draw[seg] (sec) -- node[above, font=\scriptsize]{tc} (min);
		\draw[seg] (min) -- node[above, font=\scriptsize]{tc} (ore);
	\end{tikzpicture}
	\caption{Contatori sincroni in cascata: ogni contatore abilita il successivo con il suo terminal count.}
	\label{fig:crono}
\end{figure}

\todo{Contatori binari o BCD, gestione di set e reset, valori non validi.}

\implementazione
\todo{Codice del contatore generico (in appendice) e del cronometro strutturale.}

\simulazione
\todo{Testbench con base dei tempi accelerata; passaggi 59$\rightarrow$00 e 23:59:59$\rightarrow$00:00:00.}

\subsection{Esercizio 5.2}

\begin{traccia}
	Sintetizzare ed implementare su board il componente sviluppato al punto precedente,
	utilizzando i display a 7 segmenti per la visualizzazione dell'orario (o una combinazione
	di display e led nel caso in cui i display a disposizione siano in numero inferiore a
	quello necessario), gli switch per l'immissione dell'orario iniziale e due bottoni, uno per
	il set dell'orario e uno per il reset. Si utilizzi una codifica a scelta dello studente per
	la visualizzazione dell'orario sui display (esadecimale o decimale).
\end{traccia}

\sintesiboard
\todo{Driver del display, immissione dell'orario con 16 switch (17 bit necessari), vincoli, foto.}

\subsection{Esercizio 5.3 (solo 9 CFU)}

\begin{traccia}
	Estendere il componente sviluppato ai punti precedenti in modo che sia in grado di
	acquisire e memorizzare internamente fino ad un numero N di intertempi in corrispondenza di
	un ingresso di stop. Opzionalmente, il componente può prevedere una modalità di
	visualizzazione in cui, alla pressione di un bottone, vengano visualizzati sui display gli
	intertempi memorizzati (uno per ogni pressione).
\end{traccia}

\todo{Memoria degli intertempi, modalità di visualizzazione.}

% ----------------------------------------------------------------------------------------
\section{Esercizio 6: Sistema di lettura-elaborazione-scrittura PO\_PC}
% ----------------------------------------------------------------------------------------

\subsection{Esercizio 6.1}

\begin{traccia}
	Progettare, implementare in VHDL e verificare mediante simulazione un sistema dotato di una
	memoria ROM di N locazioni da 8 bit ciascuna, una macchina combinatoria M in grado di
	trasformare (secondo una funzione a scelta dello studente) la stringa di 8 bit letta dalla
	ROM in una stringa di 4 bit, e una memoria MEM di N locazioni che memorizza la stringa in
	output da M.

	Il sistema si avvia in corrispondenza di un segnale di START che viene fornito esternamente.
	Una volta avviato, tramite un'apposita unità di controllo che gestisce la tempificazione del
	sistema, viene scandita una locazione alla volta della ROM e viene scritta la corrispondente
	locazione di MEM. Gli indirizzi di memoria sono forniti da un contatore. Le memorie ROM e
	MEM hanno rispettivamente un read e un write sincrono.
\end{traccia}

\progetto

Il sistema è diviso in una \emph{parte operativa}, che contiene i componenti che trasformano
e memorizzano i dati, e una \emph{parte di controllo}, una FSM che genera i segnali di
comando nel ciclo giusto (Figura~\ref{fig:popc}).

\begin{figure}[H]
	\centering
	\begin{tikzpicture}[node distance=1.3cm]
		\node[blocco] (cnt) {Contatore\\indirizzi};
		\node[blocco, right=of cnt] (rom) {ROM\\$N\times 8$};
		\node[blocco, right=of rom] (m) {M};
		\node[blocco, right=of m] (mem) {MEM\\$N\times 4$};
		\draw[bus] (cnt) -- node[larghezza=$k$]{/} (rom);
		\draw[bus] (rom) -- node[larghezza=8]{/} (m);
		\draw[bus] (m) -- node[larghezza=4]{/} (mem);
		\draw[bus] (cnt.south) -- ++(0,-0.5) -| node[pos=0.25, above, font=\scriptsize]{\sig{addr}} (mem.south);
		\coordinate (mid) at ($(rom)!0.5!(m)$);
		\node[ctrl, above=2cm of mid, minimum width=5.5cm] (pc) {Unità di controllo (FSM)};
		\draw[seg] ([yshift=-2mm]pc.west) -| node[pos=0.75, right, font=\scriptsize]{\sig{en\_cnt}} ([xshift=3mm]cnt.north);
		\draw[seg] ([xshift=-3mm]cnt.north) |- node[pos=0.25, left, font=\scriptsize]{\sig{tc}} ([yshift=2mm]pc.west);
		\draw[seg] (pc.south -| rom.north) -- node[right, font=\scriptsize]{\sig{rd}} (rom.north);
		\draw[seg] ([yshift=-2mm]pc.east) -| node[pos=0.75, right, font=\scriptsize]{\sig{wr}} (mem.north);
		\draw[seg] ($(pc.north)+(0,0.7)$) node[above]{\sig{START}} -- (pc.north);
		\begin{scope}[on background layer]
			\node[draw=black!40, dashed, rounded corners, fit=(cnt)(rom)(m)(mem), inner sep=0.8cm,
			      label={[font=\sffamily\scriptsize, black!60]below:parte operativa}] {};
		\end{scope}
	\end{tikzpicture}
	\caption{Scomposizione del sistema in parte operativa e parte di controllo ($k = \lceil\log_2 N\rceil$).}
	\label{fig:popc}
\end{figure}

\todo{Automa dell'unità di controllo e diagramma temporale (latenza di lettura della ROM).}

\implementazione
\todo{Codice di parte operativa, unità di controllo e sistema.}

\simulazione
\todo{Verifica che a fine processo MEM[i] = M(ROM[i]) per ogni i.}

\subsection{Esercizio 6.2}

\begin{traccia}
	Sintetizzare ed implementare su board il componente sviluppato al punto precedente,
	utilizzando due bottoni per i segnali di read e reset rispettivamente e i led per la
	visualizzazione delle uscite della macchina istante per istante.
\end{traccia}

\sintesiboard
\todo{Significato del bottone di read (passo-passo o start), vincoli, prova su scheda.}
